\chapter{Lessons: How This Project Could Have Fooled Itself}\label{ch:lessons}

\section{Introduction}\label{sec:lessons-why}
\epigraph{To consult the statistician after an experiment is finished is often merely to ask him to
conduct a post mortem examination. He can perhaps say what the experiment died of.}{Ronald Fisher,
address to the Indian Statistical Congress \citeyearpar{fisher1938}}

Most of our time went into experiments that failed, and a fair share of the failures were failures of
\emph{method} rather than of ideas: a setting that was not what we believed it to be, a number
measured in the wrong place, or a conclusion that went further than its evidence. The repository keeps
a list of these mistakes in its working notes (\file{CLAUDE.md}), written so that they would not be
repeated. In this chapter we organise that list into a small number of principles. None of the
principles is new: many of them appear in accounts of the hidden costs of machine-learning
systems~\citep{sculley2015} and of the ways in which machine-learning papers claim more than their
evidence supports~\citep{lipton2019}. What justifies a chapter is that each of them cost this project real time or money,
and each comes with a concrete example that shows how easily it is violated.

The principles fall into five groups: knowing what is actually running
(Section~\ref{sec:lessons-config}), measuring the right quantity in the right place
(Section~\ref{sec:lessons-measure}), asking the question that the metric asks
(Section~\ref{sec:lessons-question}), claiming no more than the evidence supports
(Section~\ref{sec:lessons-claims}), and spending limited resources deliberately
(Section~\ref{sec:lessons-spend}). Section~\ref{sec:lessons-checklist} condenses them into a checklist.

\section{Knowing what is actually running}\label{sec:lessons-config}

\paragraph{Configurations must be read back from the run.} We believed that exp010 was running with
full-resolution mask supervision (\texttt{mask\_ratio=1}), because its validation loop ran for 45
iterations, which implied a batch size of 2. But the batch size says nothing about the resolution of
the masks. The run's own \file{args.yaml} said \texttt{mask\_ratio: 2}, because a silent fallback had
changed the very setting that the experiment was designed to test, and twelve GPU hours measured the
wrong configuration (Section~\ref{sec:yolo-exp010}). After any run starts, we should read its resolved
configuration and confirm that it is the one we intended. \citet{sculley2015} discuss the errors
that accumulate in the configuration of machine-learning systems under the name
\emph{configuration debt}.

\paragraph{State must be queried again before it is reported.} On one occasion we reported a run as
healthy ``twenty minutes in'', on the basis of a status check that was already out of date; the run had
failed after five minutes. The state of a remote job changes, and a report about it should rest on a
fresh query.

\paragraph{A library is source code.} Ultralytics silently replaced a learning rate that we had
requested, resumed from a folder that nobody had chosen, and halved a batch after an out-of-memory error
(Section~\ref{sec:yolo-lessons}). A rented machine failed inside our own training callbacks, because
they had been written against an assumed version of the library's internals and tested only against
that assumption (Section~\ref{sec:yolo-exp097}). When our code touches a library's internals, we should
read those internals in the exact version that we pin, and rehearse the real code against that version
on free hardware before it costs money.

\paragraph{Before launching anything, list what is already running.} A handoff note stated that an
experiment ``was prepared but not launched'', and a status query for one guessed name seemed to confirm
it. In fact the experiment was running under a different name, which Kaggle had derived from its
title rather than from its identifier (Appendix~\ref{app:infra}). A duplicate was launched and ran for
ten minutes before anyone noticed. A query about one guessed name is not a check for running work, and
a claim that something was ``not launched'' should never be inherited from a note without checking it.

\section{Measuring the right thing, in the right place}\label{sec:lessons-measure}

\paragraph{The rule that is deployed must be the rule that was validated.} We validated a fusion rule
with its second model run at a confidence floor of 0.05, and then prepared it for deployment with the
floor at 0.30. The deployment file had 979 rows where the validated rule gave 1,229, and only the count
of rows revealed the mistake (Section~\ref{sec:val-traps}). The deployment path should import the
validation path, and it should reproduce a validated number before it writes anything
(Section~\ref{sec:sub-replay}).

\paragraph{Measurements belong at the metric's resolution.} A refiner's gain of $+0.0118$ at the level
of crops turned out to be an artefact of resampling, because the same measurement gave $+0.0117$ with
no model at all (Section~\ref{sec:post-v2}). We should measure at the resolution, on the population and
in the units that the metric uses: native pixels, all records, and pooled counts.

\paragraph{Every number comes from a stage of the pipeline.} The statement that ``284 misses have zero
overlap'' described the predictions emitted after the confidence floor, and not the output of the
detector; in the full pool of candidates only 65 of those filaments had never been proposed
(Section~\ref{sec:recall-stages}). A statistic measured downstream of a filter cannot describe what the
filter removed.

\paragraph{Like must be compared with like.} SQ is an average over matched pairs only. Comparing it with
a per-instance IoU averaged over all instances compares a filtered population with an unfiltered one, and
it once made our masks look as good as those of a published method when nothing of the kind had been
shown (Section~\ref{sec:metric-sq-properties}).

\paragraph{Memory experiments need separate processes.} When we tested several memory configurations in
one process, two of them produced byte-identical out-of-memory errors although the batch size between
them had halved. The second configuration had never run, because \texttt{torch.cuda.empty\_cache()}
cannot free memory that an earlier trainer still references, and every later measurement was
contaminated (Section~\ref{sec:yolo-exp031}). Each memory experiment should run in its own process, and
the allocator settings must be made before anything imports PyTorch.

\paragraph{Silence is not a result.} A script written without an
\texttt{if \_\_name\_\_ == "\_\_main\_\_"} block imported cleanly and did nothing, twice, and its
empty output could easily have been read as a negative result. We should check that a script actually
ran before interpreting its output.

\section{Asking the question the metric asks}\label{sec:lessons-question}

\paragraph{The training objective must ask the metric's question.} The crop verifier was labelled
positive if \emph{any} annotator had drawn a filament, whereas the metric asks whether \emph{this}
annotator drew it. The model was right 72\% of the time on the first question and 17\% of the time on
the second (Section~\ref{sec:recall-crop}). Before choosing labels, we should write down the exact
question that the metric asks.

\paragraph{A synthetic test that encodes an assumption cannot test it.} Snapping masks to the edge of
the intensity scored an IoU of 0.9995 on a synthetic filament and destroyed real masks, because the
synthetic filament had the hard edge whose existence was under test (Section~\ref{sec:post-free}). A
sanity check confirms that code runs; only real data can test an idea, and we should always say which of
the two we are doing.

\paragraph{An ablation needs headroom.} Two ablations of instance decoders were run on annotated masks,
where the control already scored 0.9995, so no decoder could have shown a gain
(Section~\ref{sec:par-headroom}). The score of the control should be computed first; if it is already at
the ceiling, the ablation cannot answer its question.

\paragraph{Ceiling, operator and gate, in that order.} Several repair ideas were dropped because their
triggers failed, before anyone had measured whether the repair itself was worth anything
(Section~\ref{sec:post-protocol}). A failed trigger says nothing about the repair that it was meant to
activate.

\section{Claiming no more than the evidence supports}\label{sec:lessons-claims}

\paragraph{A result about one side of the confusion matrix is a result about that side.} False positives
occurred at the rate at which annotators disagree with each other, and our notes extended this finding
to the misses, where it was false: only 2.1\% of the annotated filaments were invisible to the detector
(Section~\ref{sec:recall-orphans}). Later, after the verifier failed, the notes over-corrected in the
other direction. Two reversals on one question within a single day came from letting a measurement of one
quantity stand in for another.

\paragraph{Failed formulations are not a proof of impossibility.} After five methods for correcting
boundaries had failed, our notes declared the annotated boundary unrecoverable by any method applied
after prediction (Section~\ref{sec:post-overreach}). The measurements were sound, but the generalisation
was not. The competition's host has described uncertain boundaries as part of the physical problem,
and has encouraged participants to concentrate on the quality of their predictions rather than on the
score alone; nothing in that description implies that useful improvement is impossible.

\paragraph{Agreement between annotators is not a bound.} Two annotators agree with each other at a PQ of
0.34 and a mean IoU of 0.58. These numbers describe people. They do not limit what a model can score
against any one of them, and our model exceeds the first of them comfortably
(Section~\ref{sec:metric-agreement}).

\paragraph{The primary source should be read.} The abstract of the MAGFiLO paper reports 1,593
observations, which suggested that the public release could double our training data. The release in
fact holds 958 distinct photographs, of which the competition already uses 887, and it contains the
annotations of the test set (Section~\ref{sec:data-history}). The apparent doubling amounted to about 71
photographs, and using them would have been cheating.

\paragraph{The first explanation of a failure is a hypothesis.} We blamed rented machines that stopped
within minutes on an exhausted account. The diagnostics that the machines had uploaded showed a division
by zero in our own code (Section~\ref{sec:train-run}). The evidence that a failure leaves behind should
be read before its cause is named.

\section{Spending limited resources deliberately}\label{sec:lessons-spend}

\paragraph{Numbers already measured should change what we do.} Validation said that a fusion rule
changed 9 instances out of 1,238. We submitted it anyway, and it turned out to be 99.3\% identical to an
earlier submission. When a measurement says that a change is negligible, it should stop the spending.

\paragraph{A schedule should fit the session.} A cosine schedule sized for 400 epochs, but cut off at
epoch 62, never left its peak learning rate (Section~\ref{sec:yolo-exp010}). A short schedule that
anneals fully is better than a long one that is cut off (Section~\ref{sec:train-optim}).

\paragraph{A sweep should be costed before it runs.} Caching dense $2048\times2048$ masks for three levels
of erosion needed about 27 GB of memory, and the process was killed. Masks kept as run-length encodings
cost kilobytes, and \texttt{pycocotools} computes IoU on them directly (Section~\ref{sec:masks-rle}). In
another case, a grid of 120 points took hours, whereas a coarse grid followed by a fine one found the
same optimum in minutes.

\paragraph{The lifecycle that will actually happen is the one to test.} A plan to chain twelve-hour
Kaggle sessions relied on one session reading the output of the previous one. We tested this for a
session that had \emph{completed}. A training session, however, ends by reaching the time limit, and is
then \emph{cancelled}, and the output of a cancelled session cannot be mounted by the next one.

\paragraph{A watcher must check what has already been done.} A background loop waited for GPU quota and
then launched a training kernel. When the quota reset, we launched the kernel by hand, and the loop,
which was still running, came within two minutes of launching a duplicate. A watcher that acts must check
whether its action has already happened, and not only whether the action is possible.

\section{A checklist}\label{sec:lessons-checklist}
Figure~\ref{fig:lessons-flow} condenses the path that an idea followed during the better weeks of the
project, from its first formulation to a submission. Each step can end an idea cheaply, before the next
and more expensive step is taken.

\begin{figure}[t]
\centering
\begin{tikzpicture}[
  st/.style={box=#1, text width=5.2cm, minimum height=0.62cm, font=\sffamily\scriptsize,
             execute at begin node={\hyphenpenalty=10000}},
  why/.style={note, anchor=west, text width=4.9cm, align=left}]
  \node[st=inkmuted] (a) at (0,0) {write down the metric's exact question and the expected gain};
  \node[st=vizblue] (b) at (0,-0.95) {measure the family's ceiling with an oracle};
  \node[st=vizblue] (c) at (0,-1.9) {measure the operator under a perfect gate};
  \node[st=vizaqua] (d) at (0,-2.85) {build the real trigger; check its break-even rate};
  \node[st=vizaqua] (e) at (0,-3.8) {validate at native resolution, pooled, all records};
  \node[st=vizorange] (f) at (0,-4.75) {gain larger than its noise? paired bootstrap};
  \node[st=vizorange] (g) at (0,-5.7) {replay the validated number in the deployment path};
  \node[st=inkmuted] (h) at (0,-6.65) {look at native-resolution predictions; then submit};
  \foreach \u/\v in {a/b, b/c, c/d, d/e, e/f, f/g, g/h} \draw[flow] (\u) -- (\v);
  \node[why] at (3.1,0) {Chapter~\ref{ch:metric}; Section~\ref{sec:recall-crop}};
  \node[why] at (3.1,-0.95) {too small a ceiling: stop here};
  \node[why] at (3.1,-1.9) {operator recovers little: fix it or stop};
  \node[why] at (3.1,-2.85) {Equation~\eqref{eq:post-gate}};
  \node[why] at (3.1,-3.8) {Sections~\ref{sec:val-traps} and~\ref{sec:post-v2}};
  \node[why] at (3.1,-4.75) {the leaderboard cannot see less than about 0.01};
  \node[why] at (3.1,-5.7) {Section~\ref{sec:sub-replay}};
  \node[why] at (3.1,-6.65) {the hosts judge predictions, not only scores};
\end{tikzpicture}
\caption{The path from an idea to a submission. Each step can end the idea cheaply, before the next,
more expensive step is taken. \drawnin{book/chapters/ch14\_lessons.tex}}
\label{fig:lessons-flow}
\end{figure}

Three questions summarise the whole chapter, and we found it useful to ask them before believing any
number. First, where exactly was the number measured: at which stage of the pipeline, on which
population, at which resolution, and in which units? Second, what does the control score under the same
measurement? Third, what would the number have been if the idea had been worthless? Until all three
questions have answers, we do not yet know what the number means.

\section*{Exercises}
\begin{exercise}{Finding the configuration}
Choose any training run in \file{kaggle/} or \file{experiments/}. Find the file that records the
configuration with which it actually ran, and list every value that differs from the defaults in its
source code.
\end{exercise}

\begin{exercise}{A lesson of your own}
Choose one experiment from Appendix~\ref{app:ledger} whose conclusion you find surprising. Write down the
measurement that supports it, the population and resolution at which the measurement was taken, and one
way in which the conclusion could be wrong.
\end{exercise}

\begin{exercise}{Designing a watcher}
Write pseudocode for a background job that waits for a Kaggle kernel to finish and then submits its
output. List every condition that it must check before submitting, including those that this chapter's
examples suggest.
\end{exercise}
